\DocumentMetadata{
  pdfversion=2.0,
  pdfstandard=ua-2,
  lang=en-US,
  tagging=on,
  tagging-setup={math/setup=mathml-SE,tabsorder=structure}
}
\documentclass[11pt,letterpaper]{article}
\usepackage[margin=0.68in,includefoot]{geometry}
\usepackage{newcomputermodern}
\usepackage{amsmath,amscd,mathtools}
\usepackage{tikz,adjustbox}
\providecommand{\square}{\mdlgwhtsquare}
\usepackage[hidelinks]{hyperref}
\hypersetup{
  pdftitle={Analysis PhD Exam, May 2001},
  pdfauthor={University of Florida Department of Mathematics},
  pdfsubject={Graduate mathematics examination},
  pdfdisplaydoctitle=true,
  bookmarksopen=true
}
\setcounter{secnumdepth}{0}
\setlength{\parindent}{0pt}
\setlength{\parskip}{0.28em}
\setlength{\emergencystretch}{3em}
\setlength{\leftmargini}{2.15em}
\setlength{\leftmarginii}{2.65em}
\setlength{\leftmarginiii}{3em}
\pagestyle{empty}
\raggedbottom
\allowdisplaybreaks
\NewTaggingSocketPlug{block/list/label}{examproblem}{%
  \tagstructbegin{tag=Lbl}\tagstructend%
  \tagstructbegin{tag=\LBody}%
  \tagstructbegin{tag=\ProblemHeadingTag,title-o={\ProblemHeadingText}}%
  \tagmcbegin{tag=\ProblemHeadingTag,actualtext={\ProblemHeadingText}}%
  #2%
  \tagmcend%
  \tagstructend%
}
\newenvironment{examproblems}[2]{%
  \def\ProblemHeadingTag{#2}%
  \AssignTaggingSocketPlug{block/list/label}{examproblem}%
  \begin{enumerate}%
  \setcounter{enumi}{#1}%
  \renewcommand{\labelenumi}{\textbf{\theenumi.}}%
  \setlength{\topsep}{0.35em}%
  \setlength{\partopsep}{0pt}%
  \setlength{\itemsep}{0.48em}%
  \setlength{\parsep}{0.18em}%
}{\end{enumerate}}
\newenvironment{examsubparts}{%
  \AssignTaggingSocketPlug{block/list/label}{default}%
  \begin{enumerate}%
  \setlength{\topsep}{0.18em}%
  \setlength{\partopsep}{0pt}%
  \setlength{\itemsep}{0.16em}%
  \setlength{\parsep}{0.1em}%
}{\end{enumerate}}
\newenvironment{examinstructions}{%
  \par\begingroup\itshape
}{%
  \par\endgroup\vspace{0.18em}
}
\makeatletter
\renewcommand\subsection{\@startsection{subsection}{2}{0pt}%
  {-1.25ex plus -.25ex minus -.1ex}%
  {0.55ex plus .1ex}%
  {\normalfont\large\bfseries}}
\renewcommand{\@maketitle}{%
  \newpage\null\vskip -1em%
  {\footnotesize\bfseries
    \href{https://www.ufl.edu/}{University of Florida}, \href{https://math.ufl.edu/}{Department of Mathematics}\par}%
  \vskip -0.5em%
  \tagstructbegin{tag=H1,title={Analysis PhD Exam, May 2001}}%
  \tagpdfparaOff%
  \tagmcbegin{tag=H1}%
    {\LARGE\bfseries\href{https://gma.math.ufl.edu/exams/analysis-phd/}{Analysis PhD Exam}, May 2001\par}%
  \tagmcend%
  \tagpdfparaOn%
  \tagstructend%
  \vskip 0.25em%
  \tagmcbegin{artifact}%
    \hrule height 1pt%
  \tagmcend%
  \vskip 1em%
}
\makeatother
\title{Analysis PhD Exam, May 2001}
\author{}
\date{}
\begin{document}
\maketitle
\thispagestyle{empty}
\subsection{I}
\begin{examinstructions}
State the following theorems:
\end{examinstructions}
\begin{examproblems}{0}{H3}
\def\ProblemHeadingText{Problem 1.}
\item
The Hahn–Banach theorem.
\def\ProblemHeadingText{Problem 2.}
\item
The monotone class theorem.
\def\ProblemHeadingText{Problem 3.}
\item
The monotone convergence theorem in \(L^1(\mu)\).
\def\ProblemHeadingText{Problem 4.}
\item
The Lebesgue dominated convergence theorem in \(L^p(\mu)\).
\def\ProblemHeadingText{Problem 5.}
\item
The density of step functions in \(L^p(\mu)\), \(1\leq p<\infty\), and \(L^\infty(\mu)\).
\def\ProblemHeadingText{Problem 6.}
\item
The computation of the seminorm \(\lVert f\rVert_p\).
\def\ProblemHeadingText{Problem 7.}
\item
The integrability with respect to a measure defined by density.
\def\ProblemHeadingText{Problem 8.}
\item
The Fubini theorem for \(\mu\times\nu\)-integrable functions.
\def\ProblemHeadingText{Problem 9.}
\item
The Lusin theorem.
\def\ProblemHeadingText{Problem 10.}
\item
The Riesz representation theorem.
\end{examproblems}
\subsection{II}
\begin{examinstructions}
State and prove the following theorems:
\end{examinstructions}
\begin{examproblems}{0}{H3}
\def\ProblemHeadingText{Problem 1.}
\item
The Egorov theorem.
\def\ProblemHeadingText{Problem 2.}
\item
The integral representation of the dual of \(L^1(\mu)\), \(\mu\) finite.
\end{examproblems}
\subsection{III}
\begin{examinstructions}
Solve the following problems:
\end{examinstructions}
\begin{examproblems}{0}{H3}
\def\ProblemHeadingText{Problem 1.}
\item
Prove that if~\(X\) is a normed space and~\(Y\) is a Banach space, then \(L(X,Y)\) is a Banach space.
\def\ProblemHeadingText{Problem 2.}
\item
Prove the existence of the conditional expectation for functions \(f\in L^1(\mu)\) and for functions \(f\in L_{\mathcal F}^1(\mu)\).
\def\ProblemHeadingText{Problem 3.}
\item
Prove that the set of step functions is dense in \(L^\infty(\mu)\).
\def\ProblemHeadingText{Problem 4.}
\item
Let \((X,\Sigma,\mu)\) be a measure space and \(\mathcal R\) a ring generating the σ-algebra~\(\Sigma\). Let \(g:X\to\mathbb R\) be a~\(\mu\)-measurable function. Prove that if \(\int_A g\,d\mu=0\) for every \(A\in\mathcal R\), then \(g=0\), \(\mu\)-a.e. Is this true for \(g:X\to F\), where~\(F\) is a Banach space?
\def\ProblemHeadingText{Problem 5.}
\item
Let \((X,\Sigma)\) be a measurable space and \((\mu_n)\) a sequence of real valued measures on~\(\Sigma\). Define \(\lambda:\Sigma\to\mathbb R_+\) by
\[
\lambda(A)=\sum_{n=1}^{\infty}\frac{1}{2^n}\frac{|\mu_n|(A)}{1+|\mu_n|(X)},\qquad A\in\Sigma.
\]
 Prove that~\(\lambda\) is finite, σ-additive on~\(\Sigma\), and \(\mu_n\ll\lambda\) for each~\(n\).
\end{examproblems}
\begin{examinstructions}
Note: Write complete computations and explanations. Do not use arrows or other unnecessary signs. Use words for explanations. Write complete sentences. State the requisite formulae for each theorem and problem.
\end{examinstructions}
\end{document}
